\section{Related work}
\label{sec:related}

\subsection{Knowledge management and document acquisition}
Public knowledge bases organize facts under explicit ontologies
\citep{auer-dbpedia-kb-2007,hoffart-yago2-kb-2013}, whereas enterprise content
management emphasizes governance over organizational documents
\citep{paivarinta-ecm-2005,tyrvaainen-ecm-2006}. Research on organizational
knowledge explains how tacit experience becomes shareable
\citep{nonaka-seci-1994,nonaka-seci-ba-2000}; adjacent work studies metadata
quality, tagging, human review and auditability
\citep{park-metadata-quality-2009,golder-tagging-2006,mosqueira-rey-hitl-2022,barredo-arrieta-xai-2020}. These strands motivate reliable placement
and governance, but do not by themselves connect those controls to evidence-
checked answers. Document-AI methods, including layout-aware analysis and
MinerU, recover text and structure from complex documents
\citep{zhong-publaynet-layout-2019,huang-layoutlmv3-docai-2022,mineru}. Near-
duplicate detection and provenance provide complementary safeguards during
ingestion \citep{manku-near-dup-2007,missier-prov-2013,cheney-provenance-2009}.
QDCVR treats acquisition quality and traceable indexing as prerequisites for
reliable retrieval rather than as preprocessing assumptions.

\subsection{Retrieval, organization and agentic search}
RAG grounds generated responses in retrieved passages \citep{rag-lewis}; dense
retrievers learn shared query--passage representations \citep{dpr}, while BM25
remains a strong lexical baseline \citep{bm25}. Hybrid retrieval combines
sparse and dense signals, including reciprocal-rank fusion
\citep{rrf-cormack}; embedding models such as BGE-M3 support multilingual and
multi-granularity search \citep{bge-m3}. Graph-based systems organize and
retrieve structured relationships: GraphRAG adds entity-relation graphs and
community summaries \citep{graphrag}, LightRAG uses dual-level retrieval
\citep{lightrag}, and RAPTOR builds hierarchical document representations
\citep{raptor}. QDCVR uses a narrower document--library--tag graph for
navigation and candidate expansion; the graph does not establish answer
evidence on its own.

Tool-using agents extend retrieval through iterative interaction
\citep{react-yao,toolformer}. MCP standardizes access to tool servers
\citep{mcp}, while Modular RAG and agentic-retrieval surveys describe
reconfigurable retrieval pipelines and their efficiency challenges
\citep{modular-rag,singh-agentic-rag}. DeepRead demonstrates structure-aware
navigation for agentic search \citep{deepread}; QDCVR similarly exposes a
catalog-first librarian lane, alongside dense and hybrid alternatives, under
explicit call and tool policies.

\subsection{Verification, correction and abstention}
Self-RAG and CRAG integrate evidence assessment or retrieval correction into
generation \citep{selfrag,crag}; FLARE triggers retrieval during generation
\citep{flare}. Other work studies selective prediction under domain shift,
ontology-grounded retrieval and citation verification
\citep{kamath,cyberbot,liu-verifiability,gao-citations}. QDCVR externalizes
verification as a reusable server-side operation: the gate fetches document
text, judges evidence and returns provenance to the answer layer. This design
makes the verifier inspectable, but not infallible.

The LLM-as-a-judge literature documents preference bias, imperfect agreement,
listwise instability and self-preference \citep{zheng-llmjudge,sun-urara,panickssery-selfpref}; verifier gaming is an additional concern
\citep{lightman-verify}. These findings motivate separating evidence
verification from rank ordering and testing abstention at the answer layer.

\subsection{Research gap}
\label{sec:gap}

Table~\ref{tab:positioning} compares representative strands rather than every
system in the literature. Prior work already studies selective question
answering, citation verifiability and retrieval correction; the distinction
here is not that abstention or evidence checks are absent from prior research.
The gap this study addresses is one of integration and of role separation
rather than of missing capabilities: QDCVR examines how acquisition checks,
content-informed placement, retrieval and answer-time evidence inspection
interact in one knowledge-base workflow, reports where implementation
and transport boundaries constrain that workflow, and measures separately the
two roles a verifier can play. The table therefore motivates an
integration question, not a universal claim that earlier systems lack reusable
verification or inspectable evidence.

The study treats three design commitments as questions to test: whether
content-informed placement and lifecycle checks expose collection state;
whether a retrieved document should be re-ranked by a verifier or retain its
fusion order; and whether agent workflows can report absence while separating
completed negative searches from execution failures. Its contribution is a
knowledge-base management workflow that owns the chain from parse admission to
answer citation, a per-question account of the answers that chain produces,
and a same-pool diagnostic showing that ordering by the verifier score, not
threshold filtering, is where the measurable loss sits. It proposes no new
ranking algorithm and does not claim that the underlying retrieval problem is
solved.

\begin{table*}[t]
\centering
\small
\setlength{\tabcolsep}{4pt}
\renewcommand{\arraystretch}{1.18}
\caption{Positioning. Each strand addresses part of the path from stored
document to answer. QDCVR studies a knowledge-base workflow combining content-informed
acquisition and placement, server-side checks when invoked, and answer-layer
reporting; the evaluation measures selected components and boundaries in
Section~\ref{sec:eval}.}
\label{tab:positioning}
\begin{tabular}{@{}>{\raggedright\arraybackslash}p{0.145\textwidth}>{\raggedright\arraybackslash}p{0.285\textwidth}>{\raggedright\arraybackslash}p{0.24\textwidth}>{\raggedright\arraybackslash}p{0.24\textwidth}@{}}
\toprule
Strand & Representative work & What it optimizes & What it leaves to the deployment\\
\midrule
Knowledge management & Public knowledge bases \citep{auer-dbpedia-kb-2007,hoffart-yago2-kb-2013}; enterprise content management \citep{paivarinta-ecm-2005,tyrvaainen-ecm-2006} & An explicit ontology, or governance over a curated collection & Whether an answer drawn from that collection is supported by the text it cites\\
Document AI & Layout-aware extraction \citep{zhong-publaynet-layout-2019,huang-layoutlmv3-docai-2022}; MinerU \citep{mineru} & Recovery of text, tables and formulas from difficult layouts & Whether what was recovered is complete enough to index, and where it is then stored\\
Retrieval & Dense \citep{dpr}, lexical \citep{bm25} and hybrid ranking \citep{rrf-cormack,bge-m3} & Rank quality over a corpus that is assumed indexed & Whether a top-ranked candidate contains evidence for the claim it is used to support\\
Graph-structured organization & GraphRAG \citep{graphrag}, LightRAG \citep{lightrag}, RAPTOR \citep{raptor} & Multi-hop and hierarchical relationships between documents & Evidentiary status: a graph edge is a navigation fact, not evidence\\
Self-reflective generation & Self-RAG \citep{selfrag}, CRAG \citep{crag}, FLARE \citep{flare}; judge-based verification \citep{zheng-llmjudge,lightman-verify} & Evidence assessment or retrieval correction in generation workflows & How verification and abstention interact with knowledge-base acquisition and organization in a knowledge-base workflow\\
Agentic tool use & ReAct \citep{react-yao}, MCP \citep{mcp}, DeepRead \citep{deepread} & Access to retrieval tools behind a stable interface & Whether configured retrieval budgets and failure states are exposed and checked in a particular deployment\\
\midrule
QDCVR (this work) & Sections~\ref{sec:system}--\ref{sec:lanes} & An integration of content-informed placement, indexed-state checks, an invoked server-side evidence gate and agent answer reporting & Claims beyond the measured corpora and configurations, bounded in Section~\ref{sec:threats}\\
\bottomrule
\end{tabular}
\end{table*}
